\documentclass[conference]{IEEEtran}
\IEEEoverridecommandlockouts
\usepackage{cite}
\usepackage{amsmath,amssymb,amsfonts}
\usepackage{algorithm}
\usepackage{algorithmic}
\usepackage{graphicx}
\usepackage{textcomp}
\usepackage{booktabs}
\usepackage{xcolor}
\usepackage{url}

\begin{document}

\title{FedRL-CS: Federated Multi-Agent Reinforcement Learning for\\
Energy- and Fairness-Aware Client Selection in Edge--Cloud Networks}

\author{
\IEEEauthorblockN{Author 1, Author 2, Author 3}
\IEEEauthorblockA{Department of Computer Science and Engineering\\
SRM University AP, Andhra Pradesh, India\\
Email: \{author1, author2, author3\}@srmap.edu.in}
}

\maketitle

\begin{abstract}
Almost every account of federated learning rests on one quiet decision that is
rarely examined: which devices actually train in a given round. A round cannot
use everyone. Every device that participates spends battery, burns mobile data,
and gives up compute, and in a real fleet the phones and sensors drop in and out
without warning. So a server has to pick a small subset each round, and that
choice quietly sets how accurate, how expensive, and how fair the whole system
turns out to be. Recent work hands this choice to a reinforcement-learning agent,
but almost always a single central one that watches the entire fleet. Real
deployments are not flat: devices sit under edge servers that sit under a cloud.
We take that structure seriously and put one selection agent on each edge
cluster. Our addition is to let those agents share what they have learned:
standard federated learning already shares the \emph{model}, and we also federate
the selection \emph{policy} across clusters, on a schedule that is decoupled from
model aggregation and with a tunable blend that runs from fully-local to
fully-shared. On a controlled simulation of $40$ clients in $4$ clusters, over
$100$ rounds and $20$ seeds, the federated selector matches a single centralized
RL agent on accuracy while using $28\%$ less device energy and cutting
participation inequality (Gini) by $15\%$; a paired test puts both gaps at
$p<10^{-18}$. A reward-weight ablation shows this is not the reward talking: the
advantage survives even when the energy and fairness reward terms are switched
off, so it comes from the structure of the selection itself. Swapping the
independent critics for a shared, centralized one lifts accuracy a further
$0.83$ points, the energy saving holds at $200$ clients, and the shared policy
tolerates differential-privacy noise with no measurable cost. We are careful to
report where the method does not win---on an easy task it only ties random
selection, and the accuracy edge appears only when selection is genuinely
hard---and the whole study runs on synthetic data, which is the first thing we
would change next.
\end{abstract}

\begin{IEEEkeywords}
federated learning, client selection, multi-agent reinforcement learning, edge
computing, energy efficiency, fairness, proximal policy optimization.
\end{IEEEkeywords}

\begin{figure*}[t]
\centering
\fbox{\parbox[c][3.0cm][c]{0.95\textwidth}{\centering \color{gray}\itshape
[\,Visual abstract to be inserted here\,]\\[4pt]
\normalfont\footnotesize (placeholder --- final graphic to be generated separately)}}
\caption{Visual abstract. Each edge cluster runs a reinforcement-learning agent
that selects which of its devices train in a round; the cloud both aggregates the
model and, separately, federates the selection policy across clusters. The result
is the same accuracy as a single centralized agent at lower energy and higher
fairness.}
\label{fig:abstract}
\end{figure*}

\section{Introduction}
A federated learning (FL) system trains a shared model by sending it out to
devices, letting each device improve it on its own data, and folding the returned
updates back together, so the raw data never leaves the device~\cite{mcmahan2017,
kairouz2021}. The part that gets little attention is that a round cannot include
everyone. A production fleet has thousands or millions of devices, and each one
that joins a round pays a real price: it drains its battery, it uses the user's
data plan, and it ties up its processor. On top of that the fleet is unreliable
---a phone goes into a tunnel, an app is closed, a battery dies mid-upload. So the
server picks a handful of devices each round, and that pick is not a detail. It
decides how quickly the model learns, how much energy the fleet spends, and
whether the same few strong devices carry the load while the weak ones are never
heard from~\cite{nishio2019,lai2021}.

There is no obvious rule for the pick. Always taking the newest, best-connected
phones is fast, but it bends the model toward those users and quietly exhausts
their hardware. Taking devices at random is fair but wastes rounds on ones that
drop out. Taking everybody gives the best data coverage at the worst energy cost.
Because the trade-off shifts with the data and never settles, it is the kind of
decision one can hand to a learning agent, and several groups have~\cite{wang2020}.
What they usually build, though, is a single central agent that watches the whole
fleet at once. That is not how the systems are laid out. Devices attach to edge
servers---base stations, regional gateways, a hospital's local machine---and the
edge servers attach to a cloud. The natural fit is one selection agent per edge
cluster, not one for everything.

The question we set out to answer is whether those per-cluster agents should learn
alone or together. An ordinary FL system already shares one thing between the
clusters: the model, through aggregation. What it does not share is the selection
strategy---each cluster's agent works it out in isolation. Our idea is to share
that too. Every so often the clusters send their selection policies up to the
cloud, the cloud averages them, and each cluster blends the shared version back
into its own. A lesson one cluster learns---say, to skip devices below a fifth of
their battery---then spreads to the rest, and no cluster has to show anyone its
device data to make that happen. Two design choices make this more than a slogan.
The first is that the policy is shared on its own cadence $T_{RL}$, unhooked from
how often the model is aggregated; the second is a single blend coefficient
$\alpha$ that turns the sharing from off to full.

Our contributions are the following.
\begin{enumerate}
\item We pose hierarchical client selection as a cooperative multi-agent RL
problem and add \emph{policy federation} with a decoupled cadence $T_{RL}$ and a
tunable blend $\alpha$ (Section~\ref{sec:method}).
\item On a fixed simulation with an enforced straggler model and a grounded
energy model in joules, and averaged over $20$ seeds, we show the federated
selector matches a centralized RL agent on accuracy while cutting energy by
$28\%$ and participation inequality by $15\%$, both at $p<10^{-18}$
(Section~\ref{sec:results}).
\item We show, through a reward-weight ablation, that the saving is structural
rather than rewarded-for: it holds when the energy and fairness reward terms are
zero.
\item We add a centralized critic, optional differential-privacy noise on the
shared policy, and a personalized variant, and report plainly where each helps
and where it does not.
\item We keep the negative results in view and release all code, configurations,
and per-seed logs.
\end{enumerate}

\section{Literature Survey}
\subsection{Federated Learning and the Selection Bottleneck}
The line of work starts with McMahan \emph{et al.}~\cite{mcmahan2017}, who showed
a fleet of devices can jointly fit one model by training locally and averaging the
results, with the raw data never pooled. Later surveys map out what stays hard:
statistical heterogeneity, systems constraints, and privacy~\cite{kairouz2021,
konecny2016}. The heterogeneity is usually reproduced with a Dirichlet label
partition~\cite{hsu2019}, which we use. None of this settles which devices to use
each round, and at fleet scale that is where much of the cost lives.

\subsection{Heuristic Client Selection}
FedCS picks as many clients as can finish inside a deadline~\cite{nishio2019};
Oort ranks clients by a utility that mixes how useful their data looks with how
fast their hardware is~\cite{lai2021}. These are good, carefully-built rules, and
we keep Oort as a baseline. Their limitation for us is that the rule is fixed in
advance; it does not adapt to a fleet whose conditions drift.

\subsection{Learning to Select}
A second line treats selection as a sequential decision problem and optimizes it
with reinforcement learning~\cite{wang2020}. This is the closest in spirit to our
work. The difference is that these methods learn one policy over the whole fleet,
whereas we place a policy on each edge cluster and, more importantly, let those
policies exchange what they learn through the same cloud that already aggregates
the model.

\subsection{Multi-Agent RL and Centralized Critics}
Our agents are trained with PPO~\cite{schulman2017} and Generalized Advantage
Estimation~\cite{gae2016}. When we give them a shared value function while keeping
their actors separate, we are following the MAPPO recipe~\cite{yu2022}. The twist
here is that our agents also pass policy parameters to one another through a
federation step, which binds the multi-agent learning to the FL communication
structure rather than leaving it as a separate training detail.

\subsection{Privacy}
Differential privacy in FL is normally applied to the model
updates~\cite{abadi2016}. We instead try adding noise to the shared \emph{policy},
and we are careful to call our result utility robustness rather than a formal
guarantee, leaving an $(\varepsilon,\delta)$ accounting for later.

\subsection{The Closest Work and What We Add}
As summarized in Table~\ref{tab:related}, prior work addresses the pieces
separately: FedAvg selects nothing in particular; heuristic selectors are fixed
rules; single-agent RL learns to select but without hierarchy or sharing. To our
knowledge, the combination of per-cluster RL selection with a federated selection
policy, on a cadence decoupled from model aggregation, has not been studied, and
that combination is what we build and measure.

\begin{table}[t]
\caption{Method-level comparison with representative prior work. The comparison is
of \emph{approaches}, not numerical results.}
\label{tab:related}
\centering
\footnotesize
\begin{tabular}{p{1.5cm}p{3.0cm}p{2.6cm}}
\toprule
Work & Selection approach & What we add \\
\midrule
FedAvg~\cite{mcmahan2017} & Random / all available; no learning. & A learned,
resource-aware selector. \\
FedCS~\cite{nishio2019}, Oort~\cite{lai2021} & Fixed heuristic (deadline / utility
score). & A policy that adapts and is shared across clusters. \\
Single-agent RL~\cite{wang2020} & One RL policy over the whole fleet. & Per-cluster
agents plus policy federation. \\
\textbf{FedRL-CS (ours)} & \textbf{Per-cluster RL actors; policy federated on a
decoupled cadence $T_{RL}$, blend $\alpha$; optional central critic.} &
\textbf{---} \\
\bottomrule
\end{tabular}
\end{table}

\section{System and Problem Model}
\subsection{Network and Cluster Model}
We consider $N$ clients split across $C$ edge clusters; cluster $c$ holds a client
set $\mathcal{C}_c$. A cloud server keeps the global model $\theta$. Each client
$k$ owns a local dataset $D_k$ drawn non-IID from a shared distribution, and
carries a system state that drifts over a run: compute speed, battery level,
upload bandwidth, latency, and a staleness counter. What travels on the network is
small---model weights down, model updates up---so the sensitive data stays on the
device.

\subsection{The Round}
In round $t$ each cluster selects at most $k_{\mathrm{sel}}$ of its available
clients. The selected clients train locally and try to upload their update; the
upload is simulated with bandwidth contention and a round deadline, and any update
that misses the deadline is dropped rather than silently kept. The arrived updates
are aggregated by FedAvg and the global model is evaluated.

\subsection{What We Measure}
We report three things every round, not just the last accuracy. The first is
global test accuracy. The second is the energy the selected clients spent. The
third is the Gini coefficient of cumulative participation, which is $0$ when every
client has been used equally often and approaches $1$ when a few devices do all
the work. Because these three pull against each other, we do not assume a single
best method and instead compare on all three.

\section{Proposed Framework: FedRL-CS}
\label{sec:method}
\subsection{Overview}
Figure~\ref{fig:arch} gives the picture. Each cluster runs an RL actor that picks
clients; the cloud does two separate jobs, aggregating the model every $T_{FL}$
rounds and federating the selection policies every $T_{RL}$ rounds. We describe
each part below.

\begin{figure}[t]
\centering
\includegraphics[width=0.92\columnwidth]{figs/architecture.png}
\caption{FedRL-CS. Per-cluster RL actors select clients; the cloud aggregates the
model every $T_{FL}$ rounds and, on a separate cadence $T_{RL}$, federates the
selection policies with blend $\alpha$. The centralized critic (MAPPO variant)
supplies a shared value baseline.}
\label{fig:arch}
\end{figure}

\subsection{Learning to Select}
For cluster $c$ the actor sees a state matrix $s_c\in\mathbb{R}^{N_c\times5}$, one
row per client holding compute, battery, bandwidth, staleness, and last local
loss, together with an availability mask, and it picks $k_{\mathrm{sel}}$ clients
one at a time. After the round it receives a scalar reward
\begin{equation}
r = w_{a}\,\Delta\mathrm{acc} - w_{e}\,E - w_{f}\,\mathrm{Gini},
\label{eq:reward}
\end{equation}
where $\Delta\mathrm{acc}$ is the round's accuracy gain, $E$ the energy of the
selected clients, and $\mathrm{Gini}$ the participation inequality. Using the
accuracy \emph{delta} rather than the absolute value keeps the reward from drifting
as the model improves. Each actor is a small shared-encoder network trained with
PPO, and we hold the PPO update cadence identical across every RL method so that
none of them gets a larger gradient budget than another.

\subsection{Federating the Policy}
The new part is a second sharing channel. Every $T_{RL}$ rounds the cloud averages
the cluster policies into a meta-policy $\bar{\phi}=\frac{1}{C}\sum_c\phi_c$ and
sends back a soft update,
\begin{equation}
\phi_c \leftarrow \alpha\,\phi_c + (1-\alpha)\,\bar{\phi},
\label{eq:soft}
\end{equation}
where $\alpha\!=\!1$ leaves each policy alone (no federation) and $\alpha\!=\!0$
replaces it entirely with the consensus. Because $T_{RL}$ is set independently of
$T_{FL}$, the selection strategy and the model do not have to move in lockstep;
that decoupling is the technical heart of the method. Only policy parameters cross
the network, so the privacy posture is the same as ordinary FL.

\subsection{A Shared Critic}
Independent per-cluster critics have no way to credit a selection that was good for
the fleet as a whole. So we add a variant in which the actors stay separate but
share one centralized critic $V(\bar{s})$ over a pooled global state, and all
actors are updated against the shared advantage---the MAPPO
recipe~\cite{yu2022}---while the policy federation runs on the actors as before.

\begin{algorithm}[t]
\caption{FedRL-CS round (federated multi-agent variant)}
\label{alg:main}
\begin{algorithmic}[1]
\STATE \textbf{for} each round $t=1\ldots T$ \textbf{do}
\STATE \quad refresh client availability (churn, battery)
\STATE \quad \textbf{for} each cluster $c$: $A_c \leftarrow \pi_{\phi_c}(s_c)$
\STATE \quad clients in $\bigcup_c A_c$ train locally; incur energy $E$
\STATE \quad simulate uploads; drop stragglers past the deadline
\STATE \quad aggregate arrived updates (FedAvg); evaluate accuracy
\STATE \quad compute reward \eqref{eq:reward}; store transitions
\STATE \quad \textbf{if} $t \bmod T_{\text{PPO}}=0$: PPO-update each actor
\STATE \quad \textbf{if} $t \bmod T_{RL}=0$: federate policies via \eqref{eq:soft}
\STATE \textbf{end for}
\end{algorithmic}
\end{algorithm}

\section{Experimental Setup}
\label{sec:setup}
\subsection{Data and Partition}
We run on a synthetic multi-class classification task and split it across the
clients with a Dirichlet label skew ($\alpha_D=0.5$), so no two clients see the
same class mix~\cite{hsu2019}. The model is a multilayer perceptron trained by
FedAvg. We chose synthetic tabular data deliberately, to keep the runs fast on a
CPU and to keep the study focused on selection rather than on model architecture;
the cost of that choice is discussed in Section~\ref{sec:limits}.

\subsection{Energy and Network}
Energy is reported in joules using a grounded model that adds local-training
compute to the radio cost of an upload,
\begin{equation}
E_k = e_{\text{mac}}\,\frac{P\,n_k\,e_{\text{ep}}\,f}{\eta_k} + e_{\text{bit}}\,8B,
\end{equation}
with model parameters $P$, local samples $n_k$, epochs $e_{\text{ep}}$, a
forward--backward factor $f$, device efficiency $\eta_k$, and upload size $B$
bytes. The constants are order-of-magnitude mobile values and are left as
configuration. Uploads run over a contended link with a round deadline tuned so
that $8$--$15\%$ of them straggle and are dropped.

\subsection{Baselines and Metrics}
We compare against Random, Full Participation, Oort~\cite{lai2021}, and a single
Centralized RL agent, and we report final accuracy, cumulative energy, and
final-round Gini. The comparison that matters most is FedRL-CS against Centralized
RL, because it asks directly whether federating the policy beats simply having one
large agent. Every configuration is repeated over $20$ seeds and we quote the mean
with its standard deviation; for the main contrast we add a paired $t$-test and a
$95\%$ bootstrap confidence interval.

\subsection{Implementation}
The framework is written in Python on PyTorch. Each run writes its config, its
data split, and its seed to a file, so every number in this paper can be traced
back to the settings that produced it.

\section{Results and Discussion}
\label{sec:results}
\subsection{The Main Comparison}
Table~\ref{tab:main} lists the five methods. FedRL-CS lands at an accuracy that is
statistically indistinguishable from both Centralized RL and Random, while its
energy and Gini sit well below Centralized RL, Full, and Oort. Set against
Centralized RL it uses $28.5\%$ less energy and has $14.8\%$ lower Gini. The
paired test over the $20$ seeds makes both gaps decisive---energy gap $22.8$\,J,
$95\%$\,CI $[21.9,23.8]$, $p\approx1.2\times10^{-20}$; Gini gap $0.085$, CI
$[0.080,0.089]$, $p\approx7.1\times10^{-19}$---while the accuracy gap of
$+0.69$\,pp does not reach significance ($p=0.07$). Figure~\ref{fig:tradeoff}
shows why the picture is clean: FedRL-CS sits on the low-energy edge at the same
height as the best baselines. The reason is not mysterious. Per-cluster selection
is forced to spread the training load across clusters, so it avoids leaning on the
weak, low-efficiency devices that a single global agent keeps drafting because
they happen to look useful.

\begin{table}[t]
\caption{Main comparison ($40$ clients, $4$ clusters, $100$ rounds, $20$ seeds),
mean$\pm$std. Energy in joules; lower Gini is fairer.}
\label{tab:main}
\centering
\begin{tabular}{lccc}
\toprule
Method & Accuracy (\%) & Energy (J) & Gini \\
\midrule
\textbf{FedRL-CS (ours)} & $62.17\pm2.51$ & $\mathbf{57.32}$ & $\mathbf{0.487}$ \\
Random & $62.39\pm2.06$ & $57.42$ & $0.487$ \\
Centralized RL & $61.48\pm2.52$ & $80.12$ & $0.572$ \\
Full Participation & $61.10\pm2.72$ & $84.28$ & $0.569$ \\
Oort~\cite{lai2021} & $61.36\pm2.93$ & $55.48$ & $0.563$ \\
\bottomrule
\end{tabular}
\end{table}

\begin{figure}[t]
\centering
\includegraphics[width=0.93\columnwidth]{figs/accuracy_energy_tradeoff_v2.png}
\caption{Accuracy against energy. FedRL-CS and Random hold the low-energy edge at
the same accuracy as Centralized RL and Full Participation, which cost far more.}
\label{fig:tradeoff}
\end{figure}

\subsection{Is It Just the Cadence?}
One worry is that the result is an accident of how often PPO updates. So we ran the
two RL methods again with the update cadence moved from $1$ to $3$. Nothing of
substance changes: FedRL-CS against Centralized RL is $-28.6\%$ energy, $-14.6\%$
Gini, and again a tie on accuracy ($+0.65$\,pp, not significant). The efficiency
and fairness advantages do not depend on that setting.

\subsection{Is It Just the Reward?}
A sharper worry is circularity: the agent is rewarded for low energy and fairness,
so of course it scores well on them. We test this head-on by sweeping the reward
weights, including an accuracy-only setting with the energy and fairness terms
removed entirely. Even then FedRL-CS uses about $5$\,J less energy and has about
$0.04$ lower Gini than Centralized RL---essentially the same gap as with the full
reward (Figure~\ref{fig:reward}). The advantage is therefore built into the
structure of hierarchical selection, not handed to it by the reward, which is the
honest answer to the circularity question.

\begin{figure}[t]
\centering
\includegraphics[width=0.93\columnwidth]{figs/reward_ablation_v2.png}
\caption{Reward-weight ablation. FedRL-CS keeps its energy advantage over
Centralized RL even with the energy and fairness reward terms set to zero
(\emph{accuracy\_only}); the gap is structural.}
\label{fig:reward}
\end{figure}

\subsection{When Does the RL Earn Its Keep?}
On the easy task the RL selectors only tie random selection on accuracy, which
fairly raises the question of what the policy is for. Two results answer it. The
first is that the policies plainly learn: Figure~\ref{fig:learn} shows the
per-round reward climbing over training for both agents, from $-0.17$ to $-0.08$
for Centralized RL and from $-0.04$ to $-0.02$ for FedRL-CS, with FedRL-CS ahead
throughout. The second is that when we make selection actually matter---an extreme
non-IID split ($\alpha_D\!=\!0.1$) with a tight budget ($k_{\mathrm{sel}}\!=\!3$)
---FedRL-CS reaches $46.1\%$ accuracy against $43.6\%$ for random ($+2.48$\,pp) and
$44.4\%$ for Centralized RL ($+1.73$\,pp), while still spending far less energy and
staying fairer. The RL edge shows up exactly where the problem is hard, which is
where it should.

\begin{figure}[t]
\centering
\includegraphics[width=0.93\columnwidth]{figs/learning_curve_v2.png}
\caption{RL learning curves (per-round reward, $5$-round moving average, $20$
seeds). Both agents improve over training; FedRL-CS reaches higher reward.}
\label{fig:learn}
\end{figure}

\subsection{A Shared Critic Helps a Little}
Replacing the independent critics with one centralized critic raises accuracy from
$63.03\%$ to $63.86\%$, a gain of $0.83$ points, and it tightens the variance, at
the same energy and fairness (Table~\ref{tab:ext}). The credit-assignment signal a
shared critic provides is worth having, and the energy and fairness advantage over
Centralized RL survives the change.

\subsection{Privacy and Personalization}
Adding Gaussian noise to the shared policy, up to a standard deviation of $1.0$,
leaves accuracy, energy, and Gini untouched, which says the federated selection
policy is robust to differential-privacy noise: here, privacy is close to free. We
flag one measurement caveat honestly---because the same noised policy is sent to
every cluster, our pairwise policy-similarity metric cannot see the noise, so we
report utility robustness rather than a similarity curve, and leave a formal
$(\varepsilon,\delta)$ analysis for later. A personalized variant that shares only
the backbone and keeps each cluster's decision head local behaves as designed, with
policy similarity moving from $1.00$ when fully shared to $0.57$ when personalized
to $0.32$ when fully local, but it does not change the task numbers on this
benchmark. We report it as a mechanism that works, not as a win.

\subsection{Scaling Up}
At $200$ clients in $8$ clusters the energy saving holds: FedRL-CS uses $9\%$ less
energy than Centralized RL at tied accuracy. The fairness gap narrows to $3\%$, as
more clusters dilute any single agent's ability to over-concentrate on a few
devices, but the core efficiency benefit is clearly not an artifact of small scale.

\begin{table}[t]
\caption{Extensions: centralized critic, and scale to $200$ clients. Energy in
joules.}
\label{tab:ext}
\centering
\begin{tabular}{lccc}
\toprule
Setting & Accuracy (\%) & Energy (J) & Gini \\
\midrule
MAPPO (central critic) & $63.86\pm1.89$ & $55.70$ & $0.484$ \\
FedRL-CS (indep.\ critics) & $63.03\pm2.20$ & $55.54$ & $0.483$ \\
Centralized RL & $62.53\pm3.13$ & $79.07$ & $0.571$ \\
\midrule
FedRL-CS @ $200$ clients & $46.90\pm1.63$ & $94.73$ & $0.314$ \\
Centralized RL @ $200$ & $46.82\pm1.49$ & $104.66$ & $0.322$ \\
\bottomrule
\end{tabular}
\end{table}

\subsection{What the Results Add Up To}
Read together the numbers tell one consistent story. On accuracy the method does
not beat a good baseline, and we do not claim it does. What it does, and does
significantly and repeatably, is reach that same accuracy for meaningfully less
energy and with markedly fairer participation, and it does so because of how the
selection is structured rather than because of what it is rewarded for. The
architecture is also extensible---a shared critic buys a little accuracy, the
saving carries to larger fleets, and the shared policy takes privacy noise without
flinching. The one place accuracy moves is the hard-selection setting, which is
the honest signpost toward where this should be taken next.

\section{Limitations}
\label{sec:limits}
Two limits are worth stating plainly. First, the whole study runs on synthetic
data, and the energy, battery, network, and churn behaviour are modeled rather
than measured; the percentages here should be read as evidence about mechanism,
not as field-calibrated numbers. Second, the scale is modest, $40$ to $200$
clients, and the multi-agent learning uses PPO with a shared or centralized critic
rather than a heavier credit-assignment scheme. The differential-privacy result
establishes robustness of utility, not a formal privacy bound. We treat each of
these as a stated direction, not something tucked out of sight.

\section{Conclusion and Future Work}
We presented FedRL-CS, which federates the client-\emph{selection} policy---not
only the model---across the edge clusters of a two-tier FL system, on a cadence
decoupled from model aggregation. Over $20$ seeds it matches a centralized RL agent
on accuracy while cutting device energy by $28\%$ and participation inequality by
$15\%$, and a reward ablation shows the saving is structural. A shared critic adds
a little accuracy, the saving holds at larger scale, and the policy tolerates
privacy noise. The clear next step is to move off synthetic data onto a real
federated benchmark such as the LEAF suite, with a measured device-energy model,
and from there to a formal privacy accounting and larger deployments. All code,
configurations, and per-seed logs are released to support reproduction.

\begin{thebibliography}{00}
\bibitem{mcmahan2017} B.~McMahan, E.~Moore, D.~Ramage, S.~Hampson, and
B.~A.~y~Arcas, ``Communication-efficient learning of deep networks from
decentralized data,'' in \emph{Proc. AISTATS}, 2017, pp.~1273--1282.
\bibitem{kairouz2021} P.~Kairouz \emph{et al.}, ``Advances and open problems in
federated learning,'' \emph{Found. Trends Mach. Learn.}, vol.~14, no.~1--2,
pp.~1--210, 2021.
\bibitem{konecny2016} J.~Kone\v{c}n\'y \emph{et al.}, ``Federated learning:
Strategies for improving communication efficiency,'' \emph{arXiv:1610.05492},
2016.
\bibitem{hsu2019} T.-M.~H.~Hsu, H.~Qi, and M.~Brown, ``Measuring the effects of
non-identical data distribution for federated visual classification,''
\emph{arXiv:1909.06335}, 2019.
\bibitem{nishio2019} T.~Nishio and R.~Yonetani, ``Client selection for federated
learning with heterogeneous resources in mobile edge,'' in \emph{Proc. IEEE ICC},
2019, pp.~1--7.
\bibitem{lai2021} F.~Lai, X.~Zhu, H.~V.~Madhyastha, and M.~Chowdhury, ``Oort:
Efficient federated learning via guided participant selection,'' in \emph{Proc.
USENIX OSDI}, 2021, pp.~19--35.
\bibitem{wang2020} H.~Wang, Z.~Kaplan, D.~Niu, and B.~Li, ``Optimizing federated
learning on non-IID data with reinforcement learning,'' in \emph{Proc. IEEE
INFOCOM}, 2020, pp.~1698--1707.
\bibitem{schulman2017} J.~Schulman, F.~Wolski, P.~Dhariwal, A.~Radford, and
O.~Klimov, ``Proximal policy optimization algorithms,'' \emph{arXiv:1707.06347},
2017.
\bibitem{gae2016} J.~Schulman, P.~Moritz, S.~Levine, M.~Jordan, and P.~Abbeel,
``High-dimensional continuous control using generalized advantage estimation,'' in
\emph{Proc. ICLR}, 2016.
\bibitem{yu2022} C.~Yu \emph{et al.}, ``The surprising effectiveness of PPO in
cooperative multi-agent games,'' in \emph{Proc. NeurIPS Datasets and Benchmarks},
2022.
\bibitem{abadi2016} M.~Abadi \emph{et al.}, ``Deep learning with differential
privacy,'' in \emph{Proc. ACM CCS}, 2016, pp.~308--318.
\end{thebibliography}

\end{document}
